\chapter{تشخیص بیماری‌های آلرژیک}
\label{chap:dx-allergy}
\begingroup
\providecommand{\dxen}[1]{\lr{#1}}
\providecommand{\dxmark}[1]{\textbf{#1}}
\providecommand{\dxuncertain}[1]{\textbf{[خوانش نامطمئن: #1]}}
\providecommand{\dxpage}[2]{\clearpage\section*{صفحهٔ دست‌نویس #1}\noindent\dxen{#2}\par\medskip}
\setlist[itemize]{itemsep=0.35em,topsep=0.4em}

\dxpage{۱}{20261005\_024026281.jpg}
\noindent\textbf{تشخیص بیماری‌های آلرژیک ـ ۱۸۳}\par
\noindent\dxen{Subject: Nelson 22}\par
بیماران اغلب بیش از \dxmark{یک بیماری آلرژیک} دارند.

سابقهٔ خانوادگی بیماری‌های آلرژیک یکی از مهم‌ترین ریسک‌فاکتورهاست.
\begin{itemize}
\item ۵۰\% آتوپی یک والد؛ \dxuncertain{۶۶\%} آتوپی دو والد.
\item از طرف مادری قوی‌تر است.
\end{itemize}

\subsection*{رفتارهای کاراکتریستیک}
\begin{itemize}
\item \dxmark{\dxen{Allergic salute}}: به دلیل خارش بینی و رینوره، با کف دست بینی [را می‌مالد]. \dxmark{\dxen{Rise of the nasal crease}}.
\item \dxen{\sout{Allergic clucks}}: مالیدن کنار و پشت چشم‌ها در کونژکتیویت آلرژیک.
\item \dxmark{\dxen{Allergic clucks}}: فشار زبان به روی کام برای خارش.
\end{itemize}

\begin{itemize}
\item علائم همزمان با درو: آلرژی به قارچ یا گیاه.
\item آلرژن‌های درخت: اوایل بهار.
\item آلرژن‌های \dxen{grass}: اواخر بهار، اوایل تابستان.
\item آلرژن‌های \dxen{weeds}: اواخر تابستان، اوایل پاییز.
\item اسپور قارچ: اواخر تابستان تا اولین یخبندان زمستان.
\item آلرژن‌های چندساله: \dxen{Perennial Allergens}.
\item آلرژن‌های \dxen{indoor}: مایت، \dxen{pet}، سوسک.
\end{itemize}
آسپرژیلوس و پنی‌سیلیوم صرفاً آلرژن‌های \dxen{indoor} هستند.

\dxpage{۲}{20261005\_024032593.jpg}
\noindent\textbf{۱۸۳}\par
آلرژی‌های چندساله معمولاً همراه با \dxmark{آلرژی‌های فصلی} هستند.
\begin{itemize}
\item یک \dxen{base-line} آلرژی + بدترشدن فصلی.
\end{itemize}
شیرخواران و خردسالان زمانی به یک آلرژن حساس می‌شوند که \dxmark{برخورد مکرر} با آن داشته باشند.
\begin{itemize}
\item در مورد آلرژن‌های فصلی معمولاً باید چند سال برخورد وجود داشته باشد تا آلرژی ایجاد شود.
\end{itemize}
آلرژی‌های غذایی در \dxmark{شیرخواران و خردسالان} شایع است.
\begin{itemize}
\item تظاهرات پوستی، گوارشی و ریه، کمتر تنفسی و قلبی ـ عروقی.
\item تظاهرات ازدیاد حساسیت با واسطهٔ \dxen{IgE} معمولاً ظرف \dxmark{چند دقیقه تا ۲ ساعت} پس از خوردن [رخ می‌دهد].
\item تظاهرات غیر وابسته به \dxen{IgE} بیشتر طول می‌کشند تا بروز کنند و مزمن هستند.
\end{itemize}
\begin{itemize}
\item \dxmark{گرم‌کردن هوا} آلرژن‌های حیوانی، مایت و قارچ را تحریک می‌کند.
\item \dxmark{مرطوب‌کردن هوا} مواجهه با آلرژن مایت و قارچ را بالا می‌برد.
\item فرش و موکت مخزن آلرژن‌های مایت، قارچ و حیوانات است.
\end{itemize}
\subsection*{معاینهٔ فیزیکی ـ بیماران آلرژی}
\begin{itemize}
\item در بیمار با آسم حتماً باید \dxmark{اسپیرومتری} انجام شود.
\item اگر دیسترس تنفسی وجود داشته باشد: \dxmark{پالس‌اکسی‌متری}.
\end{itemize}

\dxpage{۳}{20261005\_024036250.jpg}
\noindent\textbf{۱۸۳}\par
\begin{itemize}
\item کودک با شکایت رینیت، رینوکونژکتیویت: بررسی از نظر تنفس دهانی، حملات عطسه، \dxen{snorting}، صاف‌کردن گلو و مالیدن بینی و چشم (نشانهٔ خارش).
\item شیرخوار باید حین شیرخوردن مشاهده شود: \dxmark{\dxen{nasal obstruction}}.
\item کودک با آسم: بررسی از نظر احتقان، سرفهٔ \dxen{wet}، تاکی‌پنه، رترکشن و ویز؛ سرفه حین گریه‌کردن.
\item \dxmark{اندازه‌گیری مکرر قد} در فواصل منظم: از نظر احتمال \dxen{FTT} و سرکوب رشد در آسم شدید و مصرف \dxen{OCS}.
\item \dxen{ICS} در کودکان قبل از بلوغ: \dxmark{\dxen{1 cm} قد کوتاه‌تر}.
\end{itemize}
وزن‌گیری ناکافی در کودک با علائم مزمن تنفسی: \dxmark{بررسی \dxen{CF}}.
\begin{itemize}
\item اندازه‌گیری \dxen{BP}: از نظر \dxen{HTN} در زمینهٔ \dxen{CS}.
\end{itemize}
\noindent\dxen{Pulsus paradoxus}: در بیمار با آسم حاد.
\begin{itemize}
\item افت بیش از \dxmark{\dxen{10 mmHg SBP}} طی دم.
\item اگر بیش از \dxmark{\dxen{20 mmHg}} باشد: مطرح‌کنندهٔ انسداد متوسط تا شدید مجاری هوایی.
\item افزایش \dxen{HR} در زمینهٔ \dxen{flare} آسم و همچنین ثانویه به آلبوترول و دکونژستان.
\item \dxmark{آلرژی باعث تب نمی‌شود}: علل عفونی بررسی شوند.
\end{itemize}

\dxpage{۴}{20261005\_024043428.jpg}
\noindent\textbf{۱۸۳}\par
\begin{itemize}
\item \dxen{Allergic shiners}: رنگ کبود ـ بنفش پلک زیری در زمینهٔ احتقان عروق.
\item در ۶۰\% بیماران آلرژی و البته در ۴۰\% بیماران بدون آلرژی.
\item مطرح‌کنندهٔ آلرژی است ولی تشخیصی نیست.
\end{itemize}
\noindent\dxmark{\dxen{Dennie--Morgan folds}}: در آتوپی شایع.
\begin{itemize}
\item در کونژکتیویت آلرژیک، درگیری چشم معمولاً قرینه است.
\item در فرم‌های شدید حتی کراتوکونوس (قوز قرنیه) نیز ممکن است.
\item کودکان تحت درمان با \dxen{CS} با دوز بالا یا مدت طولانی: ریسک ایجاد \dxen{Posterior subcapsular cataracts}.
\item \dxen{OME} ـ اوتیت مدیا با افیوژن: در رینیت آلرژیک شایع است؛ اتوسکوپی پنوماتیک.
\item کاهش یا از دست رفتن حس بویایی: بررسی از نظر پولیپ بینی و سینوزیت مزمن.
\item \dxmark{پولیپ بینی در کودکان}: شک به \dxmark{\dxen{CF}}.
\end{itemize}
مخاط بینی در رینیت آلرژیک \dxen{pale to purple} است.

\dxpage{۵}{20261005\_024047183.jpg}
\noindent\textbf{۱۸۳}\par
تندرنس سینوس فرونتال و ماگزیلاری: سینوزیت حاد بررسی شود.

هایپرتروفی تونسیلار و آدنوئیدال به همراه شرح‌حال خروپف شبانه:
\begin{itemize}
\item بررسی از نظر \dxen{Sleep Apnea}.
\end{itemize}
مشاهدهٔ \dxen{posterior pharynx}: از نظر \dxen{PND} و \dxen{PPLH} (کامل انگلیسیش).
\begin{latin}
posterior pharyngeal lymphoid hyperplasia
\end{latin}
در کودک با آسم کنترل‌شده، معاینهٔ \dxen{chest} \dxmark{باید کاملاً \dxen{N/L}} باشد.
\begin{itemize}
\item رونکای (کاهش جریان هوا) و ویزینگ در ریهٔ راست: بررسی از نظر \dxen{mucus plug} و آتلکتازی \dxen{RML}.
\item ویز دمی: بیشتر اختلالات مجاری هوایی خارج توراسیک.
\item در آسم انتظار کلابینگ انگشتان نداریم: بررسی \dxen{CF}.
\end{itemize}
\dxmark{شایع‌ترین} یافتهٔ پوستی در \dxmark{کودک با آلرژی}: \dxmark{\dxen{Xerosis}} (خشکی پوست).

\subsection*{تست‌های تشخیصی در آلرژی}
\begin{itemize}
\item ائوزینوفیلی یعنی \dxmark{بالای ۵۰۰} در هر \dxen{mcL} خون محیطی.
\item شایع‌ترین یافتهٔ هماتولوژیک بیمار آلرژی.
\end{itemize}

\dxpage{۶}{20261005\_024054192.jpg}
\noindent\textbf{۱۸۳}\par
\begin{itemize}
\item تعداد ائوزینوفیل به دنبال بروز عفونت‌ها و همچنین \dxen{OCS} کاهش می‌یابد.
\item در رینیت و ائوزینوفاژیت ائوزینوفیلیک: ائوزینوفیلی نرمال.
\item \dxmark{ائوزینوفیل \dxen{$>1500$} بدون علت مشخص}: مطرح‌کنندهٔ یکی از ۲ سندرم هایپرائوزینوفیلی.
\end{itemize}
معمولاً در بیماران آلرژی سطح \dxen{IgE} بالا است.
\begin{itemize}
\item \dxmark{\dxen{$1\,\mathrm{IU}=2.4\,\mathrm{ng}$ IgE}}.
\item طی رشد طبیعی، سطح \dxen{IgE} تا ۱۰ سالگی بالا می‌رود؛ سپس تدریجی کاهش پیدا می‌کند.
\item سطح \dxen{IgE} در \dxen{AD} \dxen{$>$} آسم آلرژیک \dxen{$>$} \dxen{AR}.
\item ۵۰\% بیماران آلرژی سطح \dxen{IgE} نرمال دارند.
\item در تشخیص آلرژیک برونکوپولمونری آسپرژیلوزیس حتماً باید سطح \dxen{Total IgE} اندازه‌گیری شود.
\item \dxmark{حتماً \dxen{$>1000\,\mathrm{ng/mL}$} باید باشد}.
\end{itemize}

\subsection*{\dxen{DDX} ائوزینوفیلی}
\begin{itemize}
\item فیزیولوژیک: \dxmark{تیره‌پوستی}، ارثی، هایپراِلِمینتیشن.
\item اندوکرین: آدیسون، هایپوپیتوئیتاریسم.
\end{itemize}
به دنبال مواجهه با آلرژن، ظرف چند دقیقه هیستامین از ماست‌سل آزاد می‌شود و ظرف ۱۰--۲۰ دقیقه به حداکثر می‌رسد و ظرف ۳۰ دقیقه بعد پایان می‌یابد.

\dxpage{۷}{20261005\_024057112.jpg}
\noindent\textbf{۱۸۳}\par
تست‌های پوستی اینترادرمال \dxmark{نباید با آلرژن‌های غذایی انجام شوند}، چون ریسک آنافیلاکسی دارند.

هر تست اینترادرمال نیاز به یک کنترل مثبت (هیستامین) و کنترل منفی (سالین) دارد.
\begin{itemize}
\item برای تشخیص \dxmark{درماتوگرافیسم}، کنترل منفی باید مثبت شود.
\item اگر کنترل مثبت، منفی شود تست بی‌ارزش.
\end{itemize}
\sout{تست‌ها}
\par
داروهایی که با تست اینترادرمال (\dxen{SPT}) تداخل دارند، باید به مدت زمان \dxmark{۵ نیمه‌عمر} قبل از تست قطع شوند.
\begin{itemize}
\item آنتی‌هیستامین‌ها.
\item آدرنرژیک‌ها (افدرین، اپی‌نفرین).
\item \dxen{CS} طولانی‌مدت هم \dxen{SPT} را مختل می‌کند.
\end{itemize}
بسیاری افراد با \dxen{SPT} مثبت هیچ شواهدی از آلرژی ندارند.

افزایش بیشتر سایز واکنش \dxen{SPT} یا سطح بالاتر \dxen{IgE} به معنی احتمال بالاتر وجود علائم بالینی است ولی ارتباطی با شدت آلرژی ندارد.

در سندرم انتروکولیت \dxen{food protein-induced}، چون \dxmark{حساسیت وابسته به \dxen{IgE} نیست}، تست \dxen{SPT} مثبت نمی‌شود.

انجام یک پنل گستردهٔ آلرژی توصیه نمی‌شود و گمراه‌کننده است.

\dxpage{۸}{20261005\_024102615.jpg}
\noindent\textbf{۱۸۳}\par
\dxmark{در پوست ملتهب یا \dxen{AD} نمی‌توان \dxen{SPT} انجام داد}.

تأیید تشخیص ائوزینوفاژیت ائوزینوفیلیک نیازمند انجام اندوسکوپی است.
\endgroup
